\uuid{O23e}
\exo7id{4156}
\titre{exo7 4156}
\auteur{quercia}
\organisation{exo7}
\datecreate{2010-03-11}
\isIndication{false}
\isCorrection{true}
\chapitre{Fonction de plusieurs variables}
\sousChapitre{Dérivée partielle}
\cpge{mpsi-mp2i.fonctions-deux-variables.b, pcsi.fonctions-deux-variables.b, ptsi.fonctions-deux-variables.b, mpsi-mp2i.primitives-equations-differentielles.b, pcsi.techniques-calcul-primitives.b, ptsi.techniques-calcul-primitives.b, bcpst1.fonctions-deux-variables.c, bcpst1.fonctions-deux-variables.a}
\module{Analyse}
\niveau{L2}
\difficulte{}

\contenu{
\begin{enumerate}
    \item \question{Déterminer les fonctions $f : {\R^{+*}\times \R^{+*}} \to  \R$ vérifiant :
    $\begin{cases} \frac{\partial f}{\partial x}(x,y) = \frac {\ln x + y - 1}{x^2y} \cr
             \frac{\partial f}{\partial y}(x,y) = \frac {\ln x }{xy^2}        \cr\end{cases}$}
\reponse{$f(x,y) = -\frac {\ln x}{xy} - \frac 1x$.}
    \item \question{Application :
    Résoudre l'équation différentielle : $(x\ln x) y' + (\ln x + y - 1)y = 0$.}
\reponse{$y = \frac {\ln x}{\lambda x - 1}$.}
\end{enumerate}
}
